% !TeX root = ../main.tex
\section{Interpreters}

We've spent a lot of time on the AST as a data structure and on environments as a
semantic device. This section makes both of those concrete: an interpreter is a
program whose job is to walk an AST, holding an environment in its hand, and produce a
result. It's the most direct way to give a language an executable meaning, direct in
the sense that there's no intermediate translation step, the tree \emph{is} the
program, and running it is just recursion over that tree.

\subsection{tree-walking, the most direct interpreter}

A \textbf{tree-walking interpreter} takes the AST built during parsing (recall the
abstract syntax trees chapter: each node is a language construct, each leaf an atomic
value) and recursively
evaluates it, node by node, according to what shape the node has. There's no separate
compilation pass, no bytecode, no machine code, the tree you parsed is the same tree
you execute. This is the simplest interpreter architecture to build and reason about,
and it's usually the first one you'd write for a new language, precisely because there's
so little machinery standing between ``I parsed a program'' and ``I can run it''.

\subsection{eval / apply}

The classic organizing structure for a tree-walking interpreter, going back to the
original Lisp interpreter, is the \textbf{eval/apply} pair:
\begin{itemize}
  \item \texttt{eval(expr, env)} looks at the \emph{shape} of \texttt{expr} (is it a
    literal? a variable reference? an if? a function call?) and dispatches to the
    matching case. Evaluating a literal just returns it. Evaluating a variable looks it
    up in \texttt{env}. Evaluating a function call means: evaluate the function
    position and the argument expressions (by calling \texttt{eval} recursively on
    each), then hand the results to \texttt{apply}.
  \item \texttt{apply(fn, args)} handles the mechanics of actually invoking a function
    once you already have a function value and already-evaluated arguments: bind the
    arguments to the function's parameters in a fresh environment (extending the
    closure's captured environment, recall the closures discussion), and \texttt{eval}
    the function's body in that new environment.
\end{itemize}
\texttt{eval} recurses over the AST's shape; \texttt{apply} recurses (indirectly,
through \texttt{eval} on the body) into the function being called. Almost every
tree-walking interpreter you'll ever read, however dressed up, is this pair underneath.

\subsection{the environment, made concrete}

The environment stopped being an abstract mathematical object ($\mathit{Var} \to
\mathit{Val}$, as in the denotational semantics chapter) the moment we started writing
an interpreter: now it's an actual runtime data structure that the interpreter's code
allocates, extends, and looks values up in. Two common concrete choices:
\begin{itemize}
  \item A \textbf{chain of frames}: each scope is a small dictionary of its own
    bindings, plus a pointer to the enclosing scope's frame. Looking up a variable
    walks the chain outward until it finds a frame that binds that name (or fails, if
    it reaches the end with nothing).
  \item A \textbf{single hashmap per call}, cloned or extended from the caller's, if
    the language doesn't need the sharing/mutation behaviour that a frame chain gives
    you for closures.
\end{itemize}
The frame-chain version is what makes closures work correctly: capturing ``the current
environment'' when a function is defined just means keeping a pointer to the current
frame, and the chain behind it, alive even after the defining scope's own \texttt{eval}
call has returned.

\subsection{a worked interpreter: arithmetic, \texttt{let}, and \texttt{if}}

Here's a complete, runnable-shaped interpreter for a tiny expression language with
integer arithmetic, \texttt{let}-bindings, and \texttt{if}, written directly over an
AST built from Python tuples for compactness:

\begin{lstlisting}[language=Python]
# AST shapes:
#   ('num', n)                       -- literal
#   ('var', name)                    -- variable reference
#   ('add', e1, e2)  ('mul', e1, e2) -- arithmetic
#   ('let', name, e1, e2)            -- let name = e1 in e2
#   ('if', cond, then_e, else_e)     -- if cond then then_e else else_e

def eval(expr, env):
    tag = expr[0]

    if tag == 'num':
        _, n = expr
        return n

    elif tag == 'var':
        _, name = expr
        return env[name]                     # look up in the environment

    elif tag == 'add':
        _, e1, e2 = expr
        return eval(e1, env) + eval(e2, env)  # recursively eval both sides

    elif tag == 'mul':
        _, e1, e2 = expr
        return eval(e1, env) * eval(e2, env)

    elif tag == 'let':
        _, name, e1, e2 = expr
        v1 = eval(e1, env)                    # eval the bound expression first
        new_env = {**env, name: v1}           # extend the environment
        return eval(e2, new_env)              # eval the body in the extended env

    elif tag == 'if':
        _, cond, then_e, else_e = expr
        if eval(cond, env) != 0:              # 0 is "false", nonzero is "true"
            return eval(then_e, env)
        else:
            return eval(else_e, env)

    else:
        raise ValueError(f"unknown AST node: {tag}")
\end{lstlisting}

\begin{example}[Tracing \texttt{let x = 2 + 3 in if x then x * 10 else 0}]
As an AST:
\begin{lstlisting}[language=Python]
program = ('let', 'x', ('add', ('num', 2), ('num', 3)),
                  ('if', ('var', 'x'),
                         ('mul', ('var', 'x'), ('num', 10)),
                         ('num', 0)))
eval(program, {})
\end{lstlisting}
The sequence of recursive calls:
\begin{enumerate}
  \item \texttt{eval(('let', 'x', \ldots), \{\})} matches the \texttt{let} case.
  \item It first evaluates the bound expression: \texttt{eval(('add', ('num',2),
    ('num',3)), \{\})}, which recurses into \texttt{eval(('num',2),\{\}) = 2} and
    \texttt{eval(('num',3),\{\}) = 3}, and returns $2+3=5$.
  \item It builds \texttt{new\_env = \{'x': 5\}} and evaluates the body,
    \texttt{eval(('if', \ldots), \{'x': 5\})}.
  \item That \texttt{if} case first evaluates the condition:
    \texttt{eval(('var','x'), \{'x':5\}) = 5} (a lookup, not a recursive descent into
    subexpressions). Since $5 \ne 0$, it's truthy.
  \item So it evaluates the \emph{then}-branch: \texttt{eval(('mul', ('var','x'),
    ('num',10)), \{'x':5\})}, which recurses into \texttt{eval(('var','x'),
    \{'x':5\}) = 5} and \texttt{eval(('num',10), \{'x':5\}) = 10}, returning
    $5 \times 10 = 50$.
  \item The \emph{else}-branch, \texttt{('num', 0)}, is never evaluated at all, exactly
    because \texttt{if} only recurses into one of its two branches.
\end{enumerate}
Final result: \texttt{50}. Notice the call tree mirrors the AST's shape exactly, that's
the entire point of tree-walking: the \emph{control flow of the interpreter} is
isomorphic to the \emph{structure of the program} being interpreted.
\end{example}

\subsection{interpreter vs.\ compiler: a spectrum, not a binary}

It's tempting to draw a hard line between ``interpreted languages'' and ``compiled
languages'', but the honest picture is a spectrum of execution strategies, and a single
language implementation can sit at more than one point on it:
\begin{itemize}
  \item \textbf{Tree-walking interpretation}: evaluate the AST directly, as above.
    Simplest to build, slowest to run (every node re-dispatches on its shape, every
    single time it's evaluated, even inside a hot loop).
  \item \textbf{Bytecode interpretation}: compile the AST once into a flat,
    lower-level instruction sequence, then run a dispatch loop over \emph{that}
    instead of re-walking the tree on every execution. We'll build this out properly
    in the next section on virtual machines; the key point for now is that ``interpreter''
    doesn't have to mean ``tree-walker'', CPython's bytecode VM and the JVM are both,
    in a real sense, interpreters too.
  \item \textbf{Native compilation}: translate all the way down to the target
    machine's own instruction set, so there's no interpreter loop running at all at
    execution time, just the CPU executing native instructions directly (the subject
    of the next section but one).
\end{itemize}
Where a real implementation sits on this spectrum is usually a pragmatic choice, not an
ideological one: many production language runtimes use several of these strategies at
once (an interpreter for a fast startup path, a JIT compiler that kicks in for code
that turns out to run often), which is a strong hint that ``interpreted'' and
``compiled'' describe \emph{implementation strategies}, not properties baked into a
language itself.

\subsection{meta-circular interpreters}

A fun (and clarifying) exercise: write an interpreter for a language \emph{in that same
language}. This is a \textbf{meta-circular interpreter}, and the classic example is
Lisp's original \texttt{eval}/\texttt{apply}, written in Lisp, evaluating Lisp:

\begin{lstlisting}[language=Lisp]
(define (eval expr env)
  (cond ((self-evaluating? expr) expr)
        ((variable? expr)       (lookup expr env))
        ((if? expr)             (if (eval (if-cond expr) env)
                                     (eval (if-then expr) env)
                                     (eval (if-else expr) env)))
        ((lambda? expr)         (make-closure expr env))
        ((application? expr)    (apply (eval (operator expr) env)
                                        (map (lambda (a) (eval a env))
                                             (operands expr))))
        (else (error "unknown expression" expr))))

(define (apply fn args)
  (if (primitive? fn)
      (apply-primitive fn args)
      (eval (closure-body fn)
            (extend-env (closure-params fn) args (closure-env fn)))))
\end{lstlisting}
This is exactly the same eval/apply shape as the Python interpreter above, just written
in Lisp itself, and evaluating Lisp expressions. Why bother? Because it forces you to
be completely explicit about what your language's semantics \emph{are}: every construct
(what counts as ``self-evaluating'', how a closure captures its environment, how
application binds arguments) has to be spelled out as code, in the same notation the
language itself uses. If you can write your language's meta-circular evaluator and it
handles every construct in the language, you've demonstrated, very concretely, that you
understand the language's semantics well enough to reproduce them.

\subsection{Exercises}
\begin{itemize}
  \item Extend the Python interpreter above with a \texttt{('sub', e1, e2)} node for
    subtraction, and a \texttt{('leq', e1, e2)} node that returns \texttt{1} if
    $e_1 \le e_2$ else \texttt{0} (recall \texttt{if} treats nonzero as true). Trace it
    on \texttt{if (5 leq 3) then 1 else 2}.
  \item Trace, by hand, the full sequence of \texttt{eval} calls (as in the worked
    example) for \texttt{let x = 4 in let y = x * x in y + 1}. What is the final
    result?
  \item The \texttt{let} case above builds \texttt{new\_env = \{**env, name: v1\}}, a
    fresh dictionary copy on every \texttt{let}. What's the asymptotic cost of this
    choice for a deeply nested chain of a thousand lets? How would using a frame-chain
    representation instead change that cost?
  \item Add a minimal \texttt{('closure', params, body)} / function-call node pair to
    the interpreter (an \texttt{eval/apply} split), following the eval/apply
    organization described above. What has to be stored in a closure value for it to
    work correctly when called later, outside the scope where it was defined?
  \item Think about it: the meta-circular evaluator's \texttt{eval} function is itself
    written in a language that presumably has its own \texttt{eval}, running somewhere
    underneath. Does this lead to infinite regress, or does it bottom out somewhere?
    Where?
  \item Compare the tree-walking interpreter's dispatch (\texttt{if tag == 'add':
    \ldots elif tag == 'mul': \ldots}) to how you'd imagine a bytecode interpreter's
    dispatch loop looking, just from the description in this section. What do you
    expect to be different about the cost of dispatching on a node/instruction when
    the tree is walked directly vs.\ when it's been flattened into a linear
    instruction stream first?
\end{itemize}
